\section{Published instance families and the standard benchmark}
\label{sec:external}
The synthetic classes and our own generators are the weakest point of any such study, because they were chosen by the authors. We therefore repeated the protocol on instances that others defined, in two steps of increasing external validity: three published families of $R||\Cmax$ instances, for which we kept the frozen learners and re-chose only the validated action (Section~\ref{sec:families}); and the standard benchmark of \citet{FanjulPeyroRuiz2010}, which comes with the best known solutions of three exact methods run by third parties and lets us report the distance to them (Section~\ref{sec:fanjul}).

\subsection{Three published families of $R||\Cmax$}
\label{sec:families}
\paragraph{Design.} The families are the uncorrelated one ($p_{jk}\sim U\{10,100\}$), the one with correlated jobs and the one with correlated machines (Section~\ref{sec:instances}), with $25$ sizes each ($m\in\{3,5,10,15,20\}$, $n\in\{30,50,80,100,200\}$), $20$ tuning and $20$ test instances per class, hence $500$ test instances per family and space. Everything that depends on the generator is re-estimated on the tuning instances (the validated action per class); the learners keep the hyperparameters frozen from the synthetic batteries, so the comparison cannot favour them by tuning to the new generator, nor penalise them by the lack of it. The unit of analysis is the mean over seeds of each instance, the cells are (family, space), and the predictions $X_1$--$X_6$ were recorded before the test run.

\paragraph{Results.} The cross-fitted headroom exceeds the frozen threshold in all six cells ($0.52$ to $2.31$ points; Table~\ref{tab:families}), so the necessary condition of the rule of use is met everywhere; the question is whether it is sufficient. The answer is mixed. The validated action beats every learner in all $30$ learner--cell pairs ($0.42$--$2.64$ points) and no learner beats it beyond the margin, as in every previous battery. But Q-Learning gains over random choice beyond the margin in four of the six cells (uncorrelated, correlated jobs and correlated machines with fifteen actions, and correlated machines with five), up to $0.67$ and $0.83$ points with correlated machines, that is $53\%$ and $65\%$ of the validated gain; with five actions it loses $0.45$ points in the uncorrelated family and is level with random choice in correlated jobs ($-0.01$). Of the six registered predictions of this battery four held. The two that did not were that Q-Learning would gain in no cell (it gained in four) and that the best learner would obtain at most $30\%$ of the validated gain wherever the validated action gains (it obtained up to $65\%$). We first read this as the most favourable evidence for online learning that we found. We withdraw that reading in Section~\ref{sec:order}: the gains are an artefact of the order of the actions (the standard learner repeats the first action of its list), and an order-invariant Q-Learning loses to random choice in all six cells.

\paragraph{Stronger competitors.} We then gave the learners the same tuning effort as the validated action and added the two published designs of Section~\ref{sec:learners}. irace received, as its budget, the cost of the tuning of the validated action ($200$ and $600$ experiments per class); it was run to find (a) a static configuration (a continuous destruction fraction and, for fifteen actions, the strategy) and (b) the four hyperparameters of Q-Learning. The static configuration of irace is at parity with the validated action in every cell (differences of at most $0.07$ points; no cell where it beats the validated action beyond the margin, as predicted in $Y_1$). The irace-tuned Q-Learning does better than the default one and gains beyond the margin in five of the six cells (predicted: at most two), up to $0.82$ points (it keeps the tie rule of the standard learner, so the caveat of Section~\ref{sec:order} applies to it as well); the two published designs gain up to $0.54$ points. The validated action beats the three in all $18$ comparisons ($0.40$--$2.19$ points), as predicted ($Y_3$, $Y_4$). Of the six predictions of this battery five held. Thus, even with the tuning of irace and with designs from the literature, a static choice validated on $20$ tuning instances per class is at least level with every alternative in every cell.

% Auto-generated by analysis/make_table_journal.py
\begin{table}[t]
\caption{Gains over random choice (points of makespan, positive = better) on the three published families of $R||C_{\max}$ (generated from the description of \citet{AbladEtAl2021}; $500$ test instances per family and space, $T=1000$ iterations). $h_{\mathrm{cf}}$: cross-fitted headroom; val: validated static action; cfg: irace-configured static action; ql, ucb, pg, alns, qo: Q-Learning, UCB1, policy gradient, ALNS roulette, optimistic Q-Learning; q: irace-tuned Q-Learning; iq, fr: the two published AOS designs. Bold: gain over random choice confirmed (Holm within method) and not equivalent to zero at the margin $\pm0.05$. No column beats the validated one beyond the margin; irace's static action is at parity with it.}
\label{tab:families}
\centering
{\scriptsize\setlength{\tabcolsep}{3.5pt}
\resizebox{\textwidth}{!}{%
\begin{tabular}{@{}llrrrrrrrrrrr@{}}
\toprule
Family & Actions & $h_{\mathrm{cf}}$ & val & cfg & ql & q & ucb & pg & alns & qo & iq & fr \\
\midrule
uncorrelated & 5 & 0.88 & 1.15 & \textbf{1.09} & $-$0.45 & $-$0.05 & $-$0.02 & 0.01 & \textbf{0.11} & $-$0.11 & $-$0.12 & \textbf{0.05} \\
correlated jobs & 5 & 0.52 & 0.55 & \textbf{0.54} & $-$0.01 & \textbf{0.14} & $-$0.01 & 0.00 & \textbf{0.12} & $-$0.06 & \textbf{0.04} & \textbf{0.05} \\
correlated machines & 5 & 1.18 & 1.26 & \textbf{1.28} & \textbf{0.67} & \textbf{0.70} & 0.00 & 0.01 & \textbf{0.50} & \textbf{0.15} & $-$0.03 & \textbf{0.30} \\
uncorrelated & 15 & 2.31 & 2.59 & \textbf{2.59} & \textbf{0.46} & \textbf{0.75} & 0.09 & 0.02 & $-$0.06 & \textbf{0.42} & \textbf{0.54} & \textbf{0.41} \\
correlated jobs & 15 & 0.79 & 0.85 & \textbf{0.82} & \textbf{0.20} & \textbf{0.35} & 0.02 & $-$0.01 & $-$0.08 & \textbf{0.04} & \textbf{0.13} & \textbf{0.07} \\
correlated machines & 15 & 1.37 & 1.29 & \textbf{1.24} & \textbf{0.83} & \textbf{0.82} & \textbf{0.08} & 0.04 & \textbf{0.28} & \textbf{0.20} & $-$0.31 & \textbf{0.25} \\
\bottomrule
\end{tabular}}}
\end{table}


\subsection{The standard benchmark of Fanjul-Peyro and Ruiz}
\label{sec:fanjul}
The benchmark of \citet{FanjulPeyroRuiz2010} (Section~\ref{sec:instances}) has $1{,}400$ instances of seven types (uniform times in $[1,100]$, $[10,100]$, $[100,120]$, $[100,200]$ and $[1000,1100]$, correlated jobs and correlated machines) with $m\in\{10,20,30,40,50\}$ and $n\in\{100,200,500,1000\}$, and best known values from CPLEX, Gurobi and an exact Lagrangian method run by third parties. We used the deterministic split of the repository's manifest (per cell of type, $m$ and $n$: two instances for tuning, five for testing, three reserved): $280$ tuning and $700$ test instances, $100$ test instances per type and action space. This is the host with the largest instances of the study, up to $n=1000$ jobs on $50$ machines.

\paragraph{Design.} The protocol is the compiled one of the earlier batteries, and nothing is tuned to this benchmark except the validated action. A class is a pair (type, $n$), $28$ in all; the validated action of a class is the best fixed action on the $10$ tuning instances of the class ($2$ seeds each). Three action spaces cover the same range as before, with the destruction sizes expressed as fractions of $n$ ($1$, $2$, $5$, $10$ and $20\%$, at least two jobs) so that the grid has the same meaning at $n=100$ and $n=1000$: B5 (five sizes of random destruction), B15 (three destruction strategies $\times$ five sizes) and B30 (B15 without and with a local search after each reconstruction). The seven learners are those of Section~\ref{sec:learners} (Q-Learning, UCB1, the policy gradient, the ALNS roulette, the optimistic Q-Learning and the two published AOS designs), with the hyperparameters frozen from the earlier batteries (those of the five- and fifteen-action synthetic spaces; B30 uses the fifteen-action values). Each method runs $T=1000$ iterations with three seeds per instance, and the unit of analysis is the mean over seeds. The cells are (type, space), $21$ in all, and the seven predictions $Z_1$--$Z_7$ and the hashes of the code and of the validated actions were recorded before the test run.

\paragraph{The validated action.} It gains over random choice beyond the margin in all $21$ cells ($Z_7$; the gain ranges from $0.14$ to $5.08$ points, Table~\ref{tab:fanjul}), and it is better than every learner in all $147$ learner--cell pairs ($Z_1$: $21$ of $21$ against Q-Learning), with no learner beating it beyond the margin in any of them ($Z_2$). In B30 the validated action is a local-search action in all $28$ classes. Measured against the best known solutions of the $700$ test instances, its mean gap is $8.32\%$ in B5, $8.32\%$ in B15 and $4.89\%$ in B30, against $9.35\%$, $10.21\%$ and $7.42\%$ for random choice; the local search helps absolute quality ($Z_4$ and $Z_5$ held, $Z_4$ narrowly at $4.89\%$ against the registered $5\%$). By type the gap of the validated action in B30 runs from $0.5\%$ ($U[1000,1100]$) to $9.9\%$ ($U[10,100]$). These gaps are those of a plain IG run for a thousand iterations, not of a method designed for this benchmark, and the protocol compares arms, not absolute quality.

\paragraph{The learners.} The learners gain over random choice beyond the margin in $39$ of the $147$ learner--cell pairs, and unevenly: Q-Learning in $3$ of its $21$ cells (all in B15, with gains of $0.07$, $0.15$ and $0.39$ points), UCB1 in one, the policy gradient in none, the ALNS roulette in six (all in B30), the optimistic Q-Learning in ten, the episodic design in seven and the FRRMAB-style design in twelve. The largest fraction of the validated gain recovered by a learner is $0.68$ (the FRRMAB-style design on $U[1000,1100]$ with five actions, where the validated gain is only $0.14$ points); in every other cell it is at most $0.46$. Two of the seven registered predictions failed. $Z_3$, that Q-Learning gains beyond the margin in at least one type in each space, failed because it gains in no type in B5 or B30. $Z_6$, that the cells where the validated action gains most are those where Q-Learning gains most (a Spearman correlation of at least $0.5$ over the $21$ cells), failed with the opposite sign: $-0.46$ ($-0.75$ in B5, $+0.14$ in B15 and $-0.68$ in B30, post hoc).

\paragraph{Reading.} The headroom here is not in doubt: the validated action gains at least $0.14$ points in every cell, far above the threshold of $0.05$, and by up to $5.08$ points. Yet the standard Q-Learning gains beyond the margin in three cells of twenty-one, loses to random choice by up to $1.19$ points, and does worst where the static choice matters most (B30, $-0.94$ points on $U[100,200]$ and $-0.88$ on correlated machines). The headroom is a necessary condition and, as in the earlier hosts, not a sufficient one; in this benchmark the relation that held on the published families of Section~\ref{sec:families} (the learner gains most where the headroom is largest) is not reproduced. We did not investigate the mechanism. One post hoc observation: in B30 the validated action evaluates $2.03$ times as many solutions per run as random choice, and Q-Learning $0.72$ times, so it concentrates on the cheaper actions while the validated action is a local-search action. The registered predictions of this battery held in five of seven cases.

% Auto-generated by analysis/make_table_journal.py
\begin{table}[t]
\caption{Standard benchmark of Fanjul-Peyro and Ruiz: gain over random choice (points of makespan, positive = better; $100$ test instances per type and space, $T=1000$ iterations, three seeds) of the validated action (val), of Q-Learning (ql) and of the best of the seven learners in each cell (best; the learner is named), and the mean gap of the validated action to the best known solution (\%). Bold: gain over random choice confirmed (Holm within learner, $21$ cells) and not equivalent to zero at the margin $\pm0.05$. B5: five destruction sizes; B15: three strategies $\times$ five sizes; B30: B15 without and with local search.}
\label{tab:fanjul}
\centering
{\scriptsize\setlength{\tabcolsep}{3pt}
\resizebox{\textwidth}{!}{%
\begin{tabular}{@{}lrrrlrrrrlrrrrlr@{}}
\toprule
 & \multicolumn{5}{c}{B5} & \multicolumn{5}{c}{B15} & \multicolumn{5}{c}{B30} \\
\cmidrule(lr){2-6}\cmidrule(lr){7-11}\cmidrule(lr){12-16}
Type & val & ql & best & & gap & val & ql & best & & gap & val & ql & best & & gap \\
\midrule
$U[1,100]$ & 2.10 & $-$1.19 & 0.42 & qo & 10.4 & 4.22 & $-$0.37 & \textbf{0.54} & fr & 10.3 & 5.08 & $-$0.57 & \textbf{0.78} & qo & 6.8 \\
$U[10,100]$ & 1.61 & $-$0.41 & 0.18 & pg & 14.9 & 3.27 & \textbf{0.39} & \textbf{0.66} & iq & 14.8 & 4.15 & $-$0.42 & \textbf{0.61} & qo & 9.9 \\
$U[100,120]$ & 0.27 & $-$0.15 & \textbf{0.12} & fr & 3.6 & 0.60 & 0.06 & \textbf{0.25} & iq & 3.6 & 0.51 & $-$0.17 & 0.10 & qo & 2.1 \\
$U[100,200]$ & 0.84 & $-$0.27 & \textbf{0.38} & qo & 11.3 & 1.56 & \textbf{0.15} & \textbf{0.32} & iq & 11.2 & 2.93 & $-$0.94 & \textbf{0.39} & alns & 4.8 \\
$U[1000,1100]$ & 0.14 & $-$0.03 & \textbf{0.10} & fr & 1.7 & 0.43 & \textbf{0.07} & \textbf{0.14} & fr & 1.7 & 0.40 & $-$0.06 & \textbf{0.09} & qo & 0.5 \\
correlated jobs & 0.24 & $-$0.16 & 0.02 & fr & 4.3 & 0.44 & $-$0.04 & 0.02 & fr & 4.3 & 0.50 & $-$0.20 & \textbf{0.13} & alns & 3.3 \\
correlated machines & 1.09 & $-$0.07 & \textbf{0.29} & fr & 12.0 & 0.94 & 0.08 & \textbf{0.26} & fr & 12.3 & 2.22 & $-$0.88 & \textbf{0.63} & alns & 6.9 \\
\bottomrule
\end{tabular}}}
\end{table}


